\section{Asymptotic behaviour of inner products}
\label{sec:inner}

In this section, we investigate asymptotic behaviour of the inner products of HDLSS
data under the spiked model. In the existing theory, the inner products have
deterministic limits. On the other hand, under the spiked model, they have random
limits. This is the source of all the results in this paper.

We write $v_{ij}=\sum_{s\le m_i}\sqrt{\lambda_{is}}z_{ijs}h_{is}$ for the spiked
part and $\xi_{ij}=\sum_{s>m_i}\sqrt{\lambda_{is}}z_{ijs}h_{is}$ for the
non-spiked part, so that $x_{ij}-\mu_i=v_{ij}+\xi_{ij}$.

\begin{lemma}
\label{lem:1}
Assume \textup{(A1)}, \textup{(A2)} and \textup{(S)}. Then, for $j\neq k$, it holds
that
\begin{equation}
\frac{(x_{ij}-\mu_i)^T(x_{ik}-\mu_i)}{\theta}
=\sum_{s=1}^{m_i}\frac{\lambda_{is}}{\theta}z_{ijs}z_{iks}+o_P(1).
\label{eq:3.1}
\end{equation}
\begin{equation}
\frac{\|x_{ij}-\mu_i\|^2}{\theta}
=\sum_{s=1}^{m_i}\frac{\lambda_{is}}{\theta}z_{ijs}^2+\frac{\tr(\Sigma_{i*})}{\theta}+o_P(1).
\label{eq:3.2}
\end{equation}
\end{lemma}

\begin{proof}
From~\eqref{eq:2.1} and $h_{is}^T h_{it}=\delta_{st}$, we have
\[
(x_{ij}-\mu_i)^T(x_{ik}-\mu_i)=\sum_{s=1}^d\lambda_{is}z_{ijs}z_{iks}.
\]
We divide the sum into $s\le m_i$ and $s>m_i$, and write
$R=\sum_{s>m_i}\lambda_{is}z_{ijs}z_{iks}$. Since $z_{ijs}$ and $z_{iks}$ are
independent for $j\neq k$, we have $E(R)=0$, so that
$\mathrm{Var}(R)=E(R^2)$. From~(A1) it holds that
\[
\mathrm{Var}(R)=\tr(\Sigma_{i*}^2).
\]
Then, from Chebyshev's inequality, for any $\varepsilon>0$ it holds that
\[
P\bigl(|R|/\theta>\varepsilon\bigr)
\le\frac{\tr(\Sigma_{i*}^2)}{\varepsilon^2\theta^2}\to 0
\]
from~(S). Thus we obtain~\eqref{eq:3.1}.

Next, we have $\|x_{ij}-\mu_i\|^2=\sum_{s=1}^d\lambda_{is}z_{ijs}^2$. Let
$Q=\sum_{s>m_i}\lambda_{is}z_{ijs}^2$. Then $E(Q)=\tr(\Sigma_{i*})$. From~(A1) we
have $E(z_{ijs}^4)\le M$ for $s>m_i$ and $E(z_{ijs}^2 z_{ijt}^2)=1$ for $s\neq t$,
so that
\[
\mathrm{Var}(Q)\le(M-1)\tr(\Sigma_{i*}^2)=o(\theta^2).
\]
Thus we obtain~\eqref{eq:3.2} from Chebyshev's inequality.
\end{proof}

\begin{lemma}
\label{lem:2}
Assume \textup{(A1)}, \textup{(A2)} and \textup{(S)}. Then, for any $j$ and $k$,
it holds that
\[
\frac{(x_{1j}-\mu_1)^T(x_{2k}-\mu_2)}{\theta}
=\sum_{s=1}^{m_1}\sum_{t=1}^{m_2}\frac{\sqrt{\lambda_{1s}\lambda_{2t}}}{\theta}\rho_{st}\,z_{1js}z_{2kt}+o_P(1).
\]
\end{lemma}

\begin{proof}
We have
\[
(x_{1j}-\mu_1)^T(x_{2k}-\mu_2)
=v_{1j}^Tv_{2k}+v_{1j}^T\xi_{2k}+\xi_{1j}^Tv_{2k}+\xi_{1j}^T\xi_{2k}.
\]
The first term is the finite sum in the statement. We evaluate the other terms.

As for $v_{1j}^T\xi_{2k}$, by conditioning on $v_{1j}$ we have
$E(v_{1j}^T\xi_{2k}\mid v_{1j})=0$ and
\[
\mathrm{Var}(v_{1j}^T\xi_{2k}\mid v_{1j})=v_{1j}^T\Sigma_{2*}v_{1j}.
\]
Here $v_{1j}^T\Sigma_{2*}v_{1j}=O_P(\theta\cdot\tr(\Sigma_{2*})/d\cdot\theta)$ is
controlled because $m_1$ is fixed and $\tr(\Sigma_{2*}^2)/\theta^2\to 0$. Since
$\|v_{1j}\|^2=O_P(\theta)$, the conditional variance is $o_P(\theta^2)$, so that
$v_{1j}^T\xi_{2k}=o_P(\theta)$. The term $\xi_{1j}^Tv_{2k}$ is handled in the same
way.

As for $\xi_{1j}^T\xi_{2k}$, from the independence we have
$E(\xi_{1j}^T\xi_{2k})=0$ and
\begin{equation}
E[(\xi_{1j}^T\xi_{2k})^2]=\tr(\Sigma_{1*}\Sigma_{2*})
\le\sqrt{\tr(\Sigma_{1*}^2)\tr(\Sigma_{2*}^2)}=o(\theta^2),
\label{eq:3.3}
\end{equation}
where we used the Cauchy--Schwarz inequality for the trace of the product of
non-negative definite matrices. Thus we obtain the result.
\end{proof}

\begin{lemma}
\label{lem:3}
Assume \textup{(A1)}, \textup{(A2)}, \textup{(S)} and \textup{(A3)} with
$\Delta>0$. Then it holds that
\[
\frac{\mu^T(x_{ij}-\mu_i)}{\theta}
=\frac{\sqrt{\Delta}}{\theta}\sum_{s=1}^{m_i}\sqrt{\lambda_{is}}\,b_{is}z_{ijs}+o_P(1)
\longrightarrow
\sqrt{\delta}\sum_{s=1}^{m_i}\sqrt{c_{is}}\,\beta_{is}z_{ijs}.
\]
\end{lemma}

\begin{proof}
From the definition of $b_{is}$ we have
$\mu^T(x_{ij}-\mu_i)=\sqrt{\Delta}\sum_{s=1}^d\sqrt{\lambda_{is}}\,b_{is}z_{ijs}$.
Let $R=\sqrt{\Delta}\sum_{s>m_i}\sqrt{\lambda_{is}}\,b_{is}z_{ijs}$. Then
$E(R)=0$ and
\[
\mathrm{Var}(R)=\Delta\sum_{s>m_i}\lambda_{is}b_{is}^2
\le\Delta\cdot\frac{\tr(\Sigma_{i*}^2)^{1/2}}{\theta}\cdot\theta
\cdot o(1),
\]
where we used $\sum_{s>m_i}b_{is}^2\lambda_{is}\le\|\mu\|^2\cdot o(1)$ from~(A3)
together with $\lambda_{i,m_i+1}/\theta\to 0$. Thus $R/\theta=o_P(1)$. The
convergence of the leading term follows from $\Delta/\theta\to\delta$,
$\lambda_{is}/\theta\to c_{is}$ and $b_{is}\to\beta_{is}$.
\end{proof}

\begin{remark}
\label{rem:1}
Lemma~\ref{lem:1} shows the essential difference from the existing theory. Under
\eqref{eq:1.2}, one has $m_i=0$ in effect and the right-hand side of~\eqref{eq:3.1}
vanishes, so that the inner product of two distinct observations converges to a
constant. Under~(S), however, the limit $\sum_{s\le m_i}c_{is}z_{ijs}z_{iks}$ is a
non-degenerate random variable. We emphasize that the term
$\tr(\Sigma_{i*})/\theta$ in~\eqref{eq:3.2} appears only for the squared norm,
namely, only when the two observations coincide. This asymmetry
between~\eqref{eq:3.1} and~\eqref{eq:3.2} plays a crucial role in
Section~\ref{sec:disc}.
\end{remark}
